\begin{center}
{\Large \bfseries \textcolor{lyred}{第一章 函数与极限}}
\end{center}

\vspace{6pt}
{\large \bfseries \textcolor{lyred}{第1.1 节 映射与函数}}

\vspace{4pt}
{\bfseries \textcolor{lyred}{一.映射}}

{\bfseries \textcolor{lyred}{1.映射概念}}

\textcolor{lyred}{定义}.设
,
X Y 为非空集合,则X 到Y 的对应法则f 称为\textcolor{lyred}{映射},若
x
X
\forall \in 
, !y
Y
\exists 
\in 
与之对应,记为
:
f
X
Y
\to 
,其中y 称为x 的\textcolor{lyred}{像},记为
()
f x , x 称为y 的一个\textcolor{lyred}{原像}. 
称X 为\textcolor{lyred}{定义域},记为
f
D ,像的全体称为\textcolor{lyred}{值域},记为
f
R ,或
(
)
f
X . 
若
(
)
f
X
Y
=
,则称f 为\textcolor{lyred}{满射}; 
若
(
)
(
)
x
x
f x
f x
\neq 
\Rightarrow 
\neq 
,即原像唯一,则称f 为\textcolor{lyred}{单射}; 
即是单射又是满射的映射称为\textcolor{lyred}{双射},或\textcolor{lyred}{一一映射}. 
{\bfseries \textcolor{lyred}{2.逆映射与复合映射}}

\textcolor{lyred}{定义}.若
:
f
X
Y
\to 
是单射,则
f
y
R
\forall \in 
, !x
X
\exists 
\in 
,使得
()
f x
y
=
,由此得到对应关系 
1 :
f
f
R
X
-
\to 
,称为f 的\textcolor{lyred}{逆映射}. 
\textcolor{lyred}{定义}.设
:
g X
Y
\to 
,
:
f Y
Z
\to 
,
Y
Y
\subset 
,则
x
X
\forall \in 
,
()
(
)
!z
f g x
Z
\exists 
=
\in 
与之对应, 
由此得到对应关系
:
f
g X
Z
\to 
\circ 
,称为g 与f 的\textcolor{lyred}{复合映射}. 
\vspace{4pt}
{\bfseries \textcolor{lyred}{二.函数}}

{\bfseries \textcolor{lyred}{1.函数的概念}}

\textcolor{lyred}{定义}.实数集之间的映射
:
f
D \to \mathbb{R}称为\textcolor{lyred}{(实)函数},习惯上记为
()
y
f x
=
, x
D
\in 
, 
其中x 称为\textcolor{lyred}{自变量}, y 称为\textcolor{lyred}{因变量},或者\textcolor{lyred}{函数值}. 
{\bfseries \textcolor{lyred}{2.反函数与复合函数}}

\textcolor{lyred}{定义}.设函数
:
f
D \to \mathbb{R}是单射,则逆映射
1 :
f
f
R
D
-
\to 
称为f 的\textcolor{lyred}{反函数},记为 
()
x
f
y
-
=
,或者
()
y
f
x
-
=
. 
\textcolor{lyred}{定义}.设有函数
()
y
f u
=
,
()
u
g x
=
,若
g
f
R
D
\subset 
,则g 与f 的复合映射称为它们的 
\textcolor{lyred}{复合函数},其中u 称为\textcolor{lyred}{中间变量}. 
{\bfseries \textcolor{lyred}{3.函数的几种初等性质}}

{\bfseries \textcolor{lyred}{(1)有界性}}

\textcolor{lyred}{定义}.若M
\exists 
\in \mathbb{R},使得当x
D
\in 
时,均有
()
f x
M
\le 
,则称
()
f x 在D 上\textcolor{lyred}{有界}; 
否则,称
()
f x 在D 上\textcolor{lyred}{无界}. 
\textcolor{lyred}{例}.证明:
()
cos
f x
x
x
=
在(
)
0,1 上无界. 
证.令
nx
n\pi 
=
,则
(
)
n
f x
n\pi 
=
,
M
\forall 
>
,当
M
n
\pi 
>
时,
(
)
n
f x
M
>
,证毕. 
\textcolor{lyred}{例}.证明:
()
sin
f x
x
x
=
在(
)
0,+\infty 上无界. 
证.令
nx
n\pi 
=
,则
(
)
n
f x
n\pi 
=
,
M
\forall 
>
,当
M
n
\pi 
>
时,
(
)
n
f x
M
>
,证毕. 
{\bfseries \textcolor{lyred}{(2)单调性}}

\textcolor{lyred}{定义}.若在区间I 上,
(
)
(
)
x
x
f x
f x
<
\Rightarrow 
<
,则称
()
f x 在I 上\textcolor{lyred}{单调增加}; 
若
(
)
(
)
x
x
f x
f x
<
\Rightarrow 
>
,则称
()
f x 在I 上\textcolor{lyred}{单调减少}. 
\textcolor{lyred}{例}.设
sin
x
y
=
,
,
2 2
y
\pi \pi 


\in -




,单调增,有反函数
arcsin
y
x
=
,
[
]
1,1
x\in -
; 
设
cos
x
y
=
,
[
]
0,
y
\pi 
\in 
,单调减,有反函数
arccos
y
x
=
,
[
]
1,1
x\in -
; 
设
tan
x
y
=
,
,
2 2
y
\pi \pi 


\in -




,单调增,有反函数
arctan
y
x
=
,
(
)
,
x\in -\infty \infty ; 
设
cot
x
y
=
,
(
)
0,
y
\pi 
\in 
,单调减,有反函数
arccot
y
x
=
,
(
)
,
x\in -\infty \infty . 
{\bfseries \textcolor{lyred}{(3)奇偶性}}

\textcolor{lyred}{定义}.设
f
D 关于原点对称,若
(
)
()
f
x
f x
-
=-
,
f
x
D
\forall \in 
,则称
()
f x 为\textcolor{lyred}{奇函数}; 
若
(
)
()
f
x
f x
-
=
,
f
x
D
\forall \in 
,则称
()
f x 为\textcolor{lyred}{偶函数}. 
{\bfseries \textcolor{lyred}{(4)周期性}}

\textcolor{lyred}{定义}.若存在
l >
,使得
(
)
()
f x
l
f x
\pm 
=
,
f
x
D
\forall \in 
,则称
()
f x 为\textcolor{lyred}{周期函数},称l 为 
\textcolor{lyred}{周期},它不是唯一的,通常说的周期是指最小正周期. 
{\bfseries \textcolor{lyred}{4.初等函数}}

\textcolor{lyred}{定义}.由常数函数及五类基本初等函数(幂函数,指数函数,对数函数,三角函数, 
反三角函数),经过有限多次的四则运算,复合所产生的,可以用一个算式表示的 
函数,称为\textcolor{lyred}{初等函数}. 
\textcolor{lyred}{例}.
1,
sgn
0,
1,
x
x
x
x
>


=
=

-
<

,[]
x 不是初等函数. 
\textcolor{lyred}{例}.
,
,
x
x
x
x
x
x
\ge 

=
=
-
<

,是初等函数. 
{\bfseries \textcolor{lyred}{5.双曲函数}}

\textcolor{lyred}{双曲正弦}sh
x
x
e
e
x
-
-
=
;\textcolor{lyred}{双曲余弦}ch
x
x
e
e
x
-
+
=
; 
\textcolor{lyred}{双曲正切}
sh
th
ch
x
x
x
x
x
e
e
x
x
e
e
-
-
-
=
=
+
. 
\textcolor{lyred}{反双曲正弦}
(
)
arsh
ln
x
x
x
=
+
+
; 
\textcolor{lyred}{反双曲余弦}
(
)
arch
ln
x
x
x
=
+
-
; 
\textcolor{lyred}{反双曲正切}
arth
ln
x
x
x
+
=
-
. 
\vspace{6pt}
{\large \bfseries \textcolor{lyred}{第1.2 节 数列的极限}}

\vspace{4pt}
{\bfseries \textcolor{lyred}{一.数列极限的定义}}

\textcolor{lyred}{定义}.给定\{\}
nx
,若存在常数a ,使得
\varepsilon 
\forall 
>
,
N
\exists 
>
,当n
N
>
时,有
nx
a
\varepsilon 
-
<
, 
则称a 为\{\}
nx
的\textcolor{lyred}{极限},记为lim
n
n
x
a
\to \infty 
=
,或
(
)
nx
a n
\to 
\to \infty . 
若\{\}
nx
的极限存在,则称该数列\textcolor{lyred}{收敛},否则称为\textcolor{lyred}{发散}. 
\vspace{4pt}
{\bfseries \textcolor{lyred}{二.收敛极限的性质}}

{\bfseries \textcolor{lyred}{1.唯一性.若}}

nx
收敛,则它的极限唯一. 
{\bfseries \textcolor{lyred}{2.有界性.若}}

nx
收敛,则它一定有界;反之不对. 
{\bfseries \textcolor{lyred}{3.保号性.若}}

(
)
lim
n
n
x
a
\to \infty 
=
>
<
,则
N
\exists 
>
,当n
N
>
时,
(
)
nx >
<
. 
\textcolor{lyred}{推论}.若除了有限多项之外,
(
)
nx \ge 
\le 
,且lim
n
n
x
a
\to \infty 
=
,则
(
)
a \ge 
\le 
. 
{\bfseries \textcolor{lyred}{4.归并性.数列}}

nx
收敛到a \Leftrightarrow 它的所有子列均收敛到a . 
\textcolor{lyred}{注}.
lim
lim
lim
n
n
n
n
n
n
x
a
x
x
a
-
\to \infty 
\to \infty 
\to \infty 
=
\Leftrightarrow 
=
=
. 
\vspace{4pt}
{\bfseries \textcolor{lyred}{三.数列极限的例子}}

\textcolor{lyred}{例}.证明:
lim 3
n
n
n
\to \infty 
+
=
+
. 
证.
(
)
3 3
nx
n
n
n
-
=
<
<
+
,
\varepsilon 
\forall 
>
,要使
nx
\varepsilon 
-
<
,只要1
n
\varepsilon 
<
,即
n
\varepsilon 
>
,故 
取
N
\varepsilon 
\varepsilon 

=
+>


,当n
N
>
时,
nx
\varepsilon 
-
<
,证毕. 
\textcolor{lyred}{例}.当
q <时, lim
n
n
q
\to \infty 
=
. 
证.
n
nx
q
-
=
,
(
)
0,1
\varepsilon 
\forall \in 
,要使
nx
\varepsilon 
-
<
,只要
n
q
\varepsilon 
<
,即lg
lg
n
q
\varepsilon 
<
,即 
lg
lg
n
q
\varepsilon 
>
,故取
lg
lg
lg
lg
N
q
q
\varepsilon 
\varepsilon 


=
+>




,当n
N
>
时,
nx
\varepsilon 
-
<
,证毕. 
\textcolor{lyred}{注}.当
a >
时, lim
n
n
a
\to \infty 
=; lim
n
n
n
\to \infty 
=. 
\vspace{4pt}
{\bfseries \textcolor{lyred}{四.数列极限的有理运算法则}}

\textcolor{lyred}{法则1}.设lim
n
n
x
\to \infty 
=
,而\{\}
ny
是有界数列,则lim
n
n
n
x y
\to \infty 
=
. 
\textcolor{lyred}{法则2}.设lim
n
n
x
a
\to \infty 
=
, lim
n
n
y
b
\to \infty 
=
,则(1)
(
)
lim
n
n
n
x
y
a
b
\to \infty 
\pm 
=
\pm ; 
(2)
(
)
lim
n
n
n
x
y
a b
\to \infty 
\cdots 
=
\cdots ;(3)当
b \neq 
时,lim
n
n
n
x
a
y
b
\to \infty 
=
. 
\textcolor{lyred}{注}.若lim
n
n
x
\to \infty 
不存在, lim
n
n
y
\to \infty 
存在,则
(
)
lim
n
n
n
x
y
\to \infty 
\pm 
一定不存在. 
\textcolor{lyred}{推论}.设某项起
n
n
x
y
\ge 
,且\{\}
nx
,\{\}
ny
收敛,则lim
lim
n
n
n
n
x
y
\to \infty 
\to \infty 
\ge 
. 
\textcolor{lyred}{例}.(1)
(
)
(
)
sin
lim
lim
sin
n
n
n
n
n
n
\to \infty 
\to \infty 
+
=
\cdots 
+
=
. 
(2)
4sin
lim 2
5sin
n
n
n
n
n
\to \infty 
-
=
+
. 
(3)
(
)
(
)
2 5
lim
lim
3 5
n
n
n
n
n
n
n
n
\to \infty 
\to \infty 
-
-
=
=-
+
+
. 
(4)
2 1
lim
lim
n
n
n
n
n
n
n
n
n
n
\to \infty 
\to \infty 
-
+
-
+
=
=
+
-
+
-
. 
\vspace{4pt}
{\bfseries \textcolor{lyred}{五.无穷小与无穷大数列}}

\textcolor{lyred}{定义}.若lim
n
n
x
\to \infty 
=
,则称\{\}
nx
为\textcolor{lyred}{无穷小}. 
若
M
\forall 
>
,
N
\exists 
>
,当n
N
>
时,
nx
M
>
,则称\{\}
nx
为\textcolor{lyred}{无穷大},记为lim
n
n
x
\to \infty 
=\infty . 
\textcolor{lyred}{定理}.设
nx \neq 
,则
lim
lim
n
n
n
n
x
x
\to \infty 
\to \infty 
=\infty \Leftrightarrow 
=
. 
\textcolor{lyred}{注}.无穷大数列一定是无界数列,反之不一定. 
\textcolor{lyred}{补充练习 }
1.
(
)
lim 1
lim
n
n
n
k
n
k
k
\to \infty 
\to \infty 
=


+
+
+
+
=
=


+
+
+
+
+
+
+


\sum 
\Lambda 
\Lambda 
. 
2.
(
)
(
)
(
)
lim
lim
lim
n
n
n
n
n
n
n
n
n n
n n
n
n
\to \infty 
\to \infty 
\to \infty 


+
+-
+
+
-
=
=
=


+
-
+
-
+
+
\Lambda 
\Lambda 
. 
3.设
lim
n
n
an
b
n
\to \infty 


-
-
=


+


,求a ,b . 
解.原式
(
)
(
)
lim
n
a
a
a n
a
b n
b
a
b
b
n
\to \infty 
-
=
=
-
-
+
-


=
=
\Rightarrow 
\Rightarrow 


+
=
=-
+


. 
\vspace{6pt}
{\large \bfseries \textcolor{lyred}{第1.3 节 函数的极限}}

\vspace{4pt}
{\bfseries \textcolor{lyred}{一.函数极限的定义}}

{\bfseries \textcolor{lyred}{1.当}}

\textcolor{lyred}{x}
\textcolor{lyred}{x}
\textcolor{lyred}{\to }
\textcolor{lyred}{时函数的极限 }
\textcolor{lyred}{定义}.设
()
f x 在某个
(
)
U x
\circ 
内有定义,若存在常数A ,使得
\varepsilon 
\forall >
,存在
\delta >
,当 
\delta 
<
-
<
x
x
时,均有
()
f x
A
\varepsilon 
-
<
,则称A 为当
x
x
\to 
时,
()
f x 的\textcolor{lyred}{极限},记为 
()
lim
x
x f x
A
\to 
=
,或
()
(
)
f x
A x
x
\to 
\to 
. 
\textcolor{lyred}{注}.(1)这里
x
x
\to 
不包含
x
x
=
的情况,故
()
lim
x
x f x
\to 
与
()
f x 在
0x 处的取值无关. 
(2)过程“
x
x
\to 
”由两部分构成:
x
x-
\to 
,
x
x+
\to 
,由此可以定义两个\textcolor{lyred}{单侧极限 }
()
lim
x
x f x
-
\to 
,
()
lim
x
x f x
+
\to 
,称为\textcolor{lyred}{左右极限}. 
(3)左极限也记为(
)
f x-,
(
)
f x -
;右极限也记为(
)
+
f x
,
(
)
f x +
. 
\textcolor{lyred}{定理}.
()
lim
x
x f x
\to 
存在
(
)
(
)
f x
f x
\Leftrightarrow 
-
=
+
. 
\textcolor{lyred}{例}.
lim
x
x C
C
\to 
=
,
lim
x
x x
x
\to 
=
,
+
lim
x
x
\to 
=
. 
\textcolor{lyred}{例}.证明:当
x >
时,
lim
x
x
x
x
\to 
=
. 
\textcolor{lyred}{例}.证明:(1)
lim sin
sin
x
x
x
x
\to 
=
;(2)
lim cos
cos
x
x
x
x
\to 
=
. 
证.(1)由
sin
sin
2sin
cos
x
x
x
x
x
x
x
x
-
+
-
=
\le 
-
,即得; 
(2)类似地,
cos
cos
x
x
x
x
-
\le 
-
,即得,证毕. 
\textcolor{lyred}{例}.设
()
sin ,
,
x
x
f x
x
x
\pi 
\pi 
\pi 
<

=
-
>

,讨论
()
lim
x
f x
\pi 
\to 
的存在性. 
解. (
)
lim sin
x
f
x
\pi 
\pi 
-
-
\to 
=
=
, (
)
lim
x
f
x
\pi 
\pi 
\pi 
+
+
\to 
=
-
=
,故
()
lim
x
f x
\pi 
\to 
=
,存在. 
\textcolor{lyred}{例}.
limsgn
x
x
\to 
不存在,因为
lim sgn
x
x
-
\to 
=-,
lim sgn
x
x
+
\to 
=.\textcolor{lyred}{ }
类似地,
[]
lim
x
x
\to 
不存在,
lim
x
x
x
\to 
-
-不存在. 
{\bfseries \textcolor{lyred}{2.当x 时函数的极限}}

\textcolor{lyred}{定义}.若存在常数A ,使得
\varepsilon 
\forall >
,
X
\exists 
>
,当x
X
>
时,
()
f x
A
\varepsilon 
-
<
,则称A 为 
当x \to \infty 时,
()
f x 的\textcolor{lyred}{极限},记为
()
lim
x
f x
A
\to \infty 
=
,或
()
(
)
f x
A x
\to 
\to \infty . 
\textcolor{lyred}{注}.过程“ x \to \infty ”由两部分构成: x \to -\infty , x \to +\infty ,故可以定义两个\textcolor{lyred}{单向极限} 
()
lim
x
f x
\to -\infty 
,
()
lim
x
f x
\to +\infty 
. 
\textcolor{lyred}{定理}.
()
lim
x
f x
\to \infty 
存在
()
()
lim
lim
x
x
f x
f x
\to -\infty 
\to +\infty 
\Leftrightarrow 
=
. 
\textcolor{lyred}{例}.limarctan
x
x
\to \infty 
不存在,因为lim arctan
x
x
\pi 
\to -\infty 
=-
, lim arctan
x
x
\pi 
\to +\infty 
=
. 
lim
x
x
x
\to \infty 
+
不存在,因为
lim
lim
x
x
x
x
x
x
\to -\infty 
\to -\infty 
+
=
=-,
lim
x
x
x
\to +\infty 
+
=+. 
\textcolor{lyred}{几何意义}.若
()
lim
x
f x
A
\to -\infty 
=
,或
()
lim
x
f x
A
\to +\infty 
=
,或
()
lim
x
f x
A
\to \infty 
=
,则称直线y
A
=
为 
曲线
()
y
f x
=
的\textcolor{lyred}{水平渐近线}. 
\textcolor{lyred}{例}.
y =
是
y
x
=
的水平渐近线;
arctan
y
x
=
有两条水平渐近线
y
\pi 
=\pm 
. 
\vspace{4pt}
{\bfseries \textcolor{lyred}{二.函数极限的性质}}

{\bfseries \textcolor{lyred}{1.唯一性.若}}

()
lim
x
x f x
\to 
存在,则极限唯一. 
{\bfseries \textcolor{lyred}{2.局部有界性.若}}

()
lim
x
x f x
\to 
存在,则在某个
(
)
U x
\circ 
中
()
f x 有界. 
{\bfseries \textcolor{lyred}{3.局部保号性.若}}

()
(
)
lim
x
x f x
A
\to 
=
>
<
,则在某个
(
)
U x
\circ 
中
()
(
)
f x >
<
. 
\textcolor{lyred}{推论}.若在某个
(
)
U x
\circ 
内
()
(
)
f x \ge 
\le 
,且
()
lim
x
x f x
A
\to 
=
,则
(
)
A \ge 
\le 
. 
\textcolor{lyred}{例}.设
()
(
)
(
)
lim
x
x
f x
f x
a
x
x
\to 
-
=
>
-
,证明:在某个
(
)
U x
内
(
)
f x
为最小值. 
证.在某个
(
)
U x
\circ 
中
()
(
)
(
)
()
(
)
f x
f x
f x
f x
x
x
-
>
\Rightarrow 
>
-
,证毕. 
{\bfseries \textcolor{lyred}{4.Heine 定理.}}

()
(
)
lim
n
n
x
x f x
A
x
x
x
x
\to 
=
\Leftrightarrow \forall 
\to 
\neq 
,均有
(
)
lim
n
n
f x
A
\to \infty 
=
. 
\textcolor{lyred}{例}.设
()
cos
f x
x
=
,证明:
()
lim
x
f x
\to 
不存在. 
证.取
nx
n\pi 
=
,则lim
n
n
x
\to \infty 
=
,
(
)
lim
n
n
f x
\to \infty 
=,再取
(
)
ny
n
\pi 
=
-
,则lim
n
n
y
\to \infty 
=
, 
(
)
lim
n
n
f
y
\to \infty 
=-,证毕. 
\textcolor{lyred}{注}.类似地,(1)若
()
lim
x
f x
\to \infty 
存在,则极限必唯一. 
(2)若
()
lim
x
f x
\to \infty 
存在,则
X
\exists 
>
,使得
()
f x 在x
X
>
上有界. 
(3)若
()
lim
x
f x
A
\to \infty 
=
,
(
)
A >
<
,则
X
\exists 
>
,当x
X
>
时
()
(
)
f x >
<
. 
(4)
()
lim
n
x
f x
A
x
\to \infty 
=
\Leftrightarrow \forall 
\to \infty ,均有
(
)
lim
n
n
f x
A
\to \infty 
=
. 
\textcolor{lyred}{例}.设
()
f x 为周期函数,证明:若
()
lim
x
f x
\to \infty 
存在,则
()
f x 为常数. 
证.设
()
lim
x
f x
A
\to \infty 
=
,若存在
0x ,使得
(
)
f x
B
A
=
\neq 
,令
nx
x
nT
=
+
,则 
nx \to \infty ,而
(
)
n
f x
B
A
\to 
\neq 
,矛盾,故
()
f x
A
\equiv 
,证毕. 
\textcolor{lyred}{补充练习 }
1.设
()
sin
f x
x
=
,证明:
()
lim
x
f x
\to \infty 
不存在. 
证.取
nx
n
\pi 
\pi 
=
+
,
ny
n
\pi 
\pi 
=
-
,则
nx \to \infty ,
ny \to \infty ,而 
(
)
lim
n
n
f x
\to \infty 
=,
(
)
lim
n
n
f
y
\to \infty 
=-,证毕. 
\vspace{6pt}
{\large \bfseries \textcolor{lyred}{第1.4 节 无穷小与无穷大}}

\vspace{4pt}
{\bfseries \textcolor{lyred}{一.无穷小}}

\textcolor{lyred}{定义}.若
()
lim
x
f x
\to 
=
\Omega 
,则称
()
f x 为当x \to \Omega 时的\textcolor{lyred}{无穷小量}. 
\textcolor{lyred}{例}.当x \to \infty 时, 1
x 为无穷小;当
x
+
\to 
时,
x -为无穷小; 
当x \to -\infty 时,
xe 为无穷小. 
\textcolor{lyred}{定理}.
()
()
()
lim
x
f x
A
f x
A
x
\alpha 
\to 
=
\Leftrightarrow 
=
+
\Omega 
,其中
()
lim
x
x
\alpha 
\to 
=
\Omega 
. 
\vspace{4pt}
{\bfseries \textcolor{lyred}{二.无穷大}}

\textcolor{lyred}{定义}.设
()
f x 在某个
(
)
U x
\circ 
中有定义,若
M
\forall 
>
,
\delta 
\exists 
>
,当
x
x
\delta 
<
-
<
时,有 
()
f x
M
>
,则称
()
f x 为当
x
x
\to 
时的\textcolor{lyred}{无穷大},记为
()
lim
x
x f x
\to 
=\infty . 
若
M
\forall 
>
,
X
\exists 
>
,当x
X
>
时,
()
f x
M
>
,则称
()
f x 为当x \to \infty 时的\textcolor{lyred}{无穷大}, 
记为
()
lim
x
f x
\to \infty 
=\infty . 
\textcolor{lyred}{例}.
lim
x
x
\to 
=\infty ,
lim
x
x
-
\to 
=-\infty ,
lim
x
x
+
\to 
=+\infty . 
\textcolor{lyred}{例}. lim
x
x
e
\to +\infty 
=\infty ,或者lim
x
x
e
\to +\infty 
=+\infty . 
\textcolor{lyred}{例}.
lim
x
x
x
e
e
\to 
-
+
不存在:
lim
lim
u
x
u
x
e
e
-
\to -\infty 
\to 
=
=
,
lim
lim
u
x
u
x
e
e
+
\to +\infty 
\to 
=
=+\infty ,故 
lim
x
x
x
e
e
-
\to 
-
=
+
,
lim
x
x
x
e
e
+
\to 
-
=-
+
. 
\textcolor{lyred}{几何意义}.若
()
lim
x
x f x
-
\to 
=\infty ,或
()
lim
x
x f x
+
\to 
=\infty ,则称直线
x
x
=
为曲线
()
y
f x
=
的 
\textcolor{lyred}{铅直渐近线}. 
\textcolor{lyred}{例}.直线
x =是
y
x
=
-的铅直渐近线. 
\vspace{4pt}
{\bfseries \textcolor{lyred}{三.无穷大与无界量的关系}}

\textcolor{lyred}{定理}.
()
(
)
lim
n
n
x
f x
x
x
\to 
=\infty \Leftrightarrow \forall 
\to 
\neq 
\Omega 
\Omega 
\Omega ,均有
(
)
lim
n
n
f x
\to \infty 
=\infty . 
\textcolor{lyred}{例}.证明:
()
sin
f x
x
x
=
不是
x \to 
时的无穷大. 
证.取
nx
n\pi 
=
,则
nx \to 
,而
(
)
n
f x
=
,证毕. 
\textcolor{lyred}{例}.证明:
()
sin
x
f x
e
x
=
不是x \to +\infty 时的无穷大. 
证.取
nx
n\pi 
=
,则
nx \to +\infty ,而
(
)
n
f x
=
,证毕. 
\textcolor{lyred}{注}.无穷大量一定是无界量,反之不一定. 
\vspace{4pt}
{\bfseries \textcolor{lyred}{四.无穷大与无穷小的关系}}

\textcolor{lyred}{定理}.设在x \to \Omega 的过程中
()
f x \neq 
,则
()
()
lim
lim
x
x
f x
f x
\to 
\to 
=\infty \Leftrightarrow 
=
\Omega 
\Omega 
. 
\textcolor{lyred}{补充练习 }
1.求曲线
x
y
x
x
-
=
+
-
的水平与铅直渐近线. 
解.
(
)(
)
(
)(
)
x
x
y
x
x
-
+
=
-
+
,
lim
lim
x
x
x
y
x
\to -
\to -
+
=
=\infty 
+
,故有铅直渐近线
x =-; 
lim
lim
x
x
x
y
x
\to \infty 
\to \infty 
+
=
=
+
,故有水平渐近线
y =. 
\vspace{6pt}
{\large \bfseries \textcolor{lyred}{第1.5 节 极限运算法则}}

\vspace{4pt}
{\bfseries \textcolor{lyred}{一.函数和差积商的极限运算法则}}

\textcolor{lyred}{定理}.(局部)有界函数与无穷小的积为无穷小. 
(1)设在
0x 的某个去心邻域内()
g x 有界,而
()
lim
x
x f x
\to 
=
,则
()()
lim
x
x f x g x
\to 
=
. 
(2)设在某个\{
\}
:
x
x
X
\ge 
上()
g x 有界,而
()
lim
x
f x
\to \infty 
=
,则
()()
lim
x
f x g x
\to \infty 
=
. 
\textcolor{lyred}{定理}.设
()
lim
x
f x
A
\to 
=
\Omega 
,
()
lim
x
g x
B
\to 
=
\Omega 
(存在),则(1)
()
()
lim
x
f x
g x
A
B
\to 
\pm 
=
\pm 




\Omega 
; 
(2)
()
()
lim
x
f x
g x
A B
\to 
\cdots 
=
\cdots 




\Omega 
;(3)当
B \neq 
时,
()
()
lim
x
f x
A
g x
B
\to 
=
\Omega 
. 
\textcolor{lyred}{注}.若
()
lim
x
f x
\to \Omega 
不存在,
()
lim
x
g x
\to \Omega 
存在,则
()
()
lim
x
f x
g x
\to 
\pm 




\Omega 
不存在. 
\textcolor{lyred}{推论}.
()
()
()
()
1 1
lim
lim
lim
n
n
n
n
x
x
x
k f
x
k f
x
k
f
x
k
f
x
\to 
\to 
\to 
+
+
=
+
+




\Omega 
\Omega 
\Omega 
\Lambda 
\Lambda 
. 
\textcolor{lyred}{推论}.
()
()
lim
lim
n
n
x
x
f x
f x
\to 
\to 


=






\Omega 
\Omega 
. 
\textcolor{lyred}{推论}.设
()
P x 为多项式,则
()
(
)
lim
x
x P x
P x
\to 
=
. 
\textcolor{lyred}{推论}.设()
()
x
x
\varphi 
\zeta 
\ge 
,若
()
lim
x
x
a
\varphi 
\to 
=
\Omega 
,
()
lim
x
x
b
\zeta 
\to 
=
\Omega 
,则a
b
\ge 
. 
\textcolor{lyred}{例}.
lim sin
x
x
x
\to 
=
,
arctan
lim
x
x
x
\to \infty 
=
. 
\textcolor{lyred}{例}.(1)
lim
5 2
x
x
x
x
\to 
-
-
=
=-
-
+
-\cdots +
. 
(2)
lim
x
x
x
x
\to 
-
=\infty 
-
+
,因为
lim
x
x
x
x
\to 
-
+
=
=
-
-
. 
(3)
(
)(
)
(
)(
)
lim
lim
lim
x
x
x
x
x
x
x
x
x
x
x
x
x
\to 
\to 
\to 
-
-
-
+
-
=
=
=
-
-
+
-
+
. 
\textcolor{lyred}{例}.(1)
lim
lim
2 1
x
x
x
x
x
x
x
x
x
x
x
\to \infty 
\to \infty 
-
-
-
-
-
-
=
=
=
-
+
-
+
-
-
. 
(2)
lim 3
x
x
x
x
x
\to \infty 
-
+
=\infty 
-
-
. 
(3)
lim 2
x
x
x
x
x
\to \infty 
-
-
=
-
+
. 
\textcolor{lyred}{注}.
(
)
0,
lim
,
lim
,
,
m
m
m
m
m
m
n
m
n
n
n
n
x
x
n
n
n
n
m
n
a x
a
x
a
a x
a
a
b
m
n
b x
b
x
b
b x
b
m
n
-
-
-
\to \infty 
\to \infty 
-

<

+
+
+

\neq 
=
=
=

+
+
+

\infty 
>

\Lambda 
\Lambda 
. 
\textcolor{lyred}{例}.
(
)(
)
(
)
(
)
(
)(
)
lim
lim
x
x
x
x
x x
x
x
x
x
x
\to \infty 
\to \infty 
+
-
-
+
+
-
+
=
=
+
+
\Lambda 
\Lambda 
\Lambda 
. 
\textcolor{lyred}{例}.(1)设
lim
x
x
ax
b
x
x
\to 
+
+
=
-
-
,求a ,b . 
解.
(
)
lim
x
x
ax
b
a
b
\to 
+
+
=
+
+
=
,故
b
a
=-
-
,于是 
(
)(
)
(
)(
)
lim
lim
x
x
x
x
a
x
ax
a
a
x
x
x
x
\to 
\to 
-
+
+
+
-
-
+
=
=
=
-
-
-
+
,故
a =
,
b =-. 
(2)设
lim
x
x
ax
b
x
\to \infty 


-
+
+
=


+


,求a ,b . 
解.原式
(
)
(
)
lim
x
a
a
a x
a
b x
b
a
b
b
x
\to \infty 
+
=
=-
+
+
+
+
-


=
=
\Rightarrow 
\Rightarrow 


+
=
=
+


. 
\vspace{4pt}
{\bfseries \textcolor{lyred}{二.复合函数的极限运算法则}}

\textcolor{lyred}{例}.
lim
lim
u
x
u
x
e
e
-
\to -\infty 
\to 
=
=
,
lim
lim
u
x
u
x
e
e
+
\to +\infty 
\to 
=
=+\infty , 
lim arctan
lim arctan
u
x
u
x
\pi 
-
\to -\infty 
\to 
=
=-
,
lim arctan
lim arctan
u
x
u
x
\pi 
+
\to +\infty 
\to 
=
=
. 
\textcolor{lyred}{例}.
lim
lim
x
u
x
u
x
\to 
\to 
+
=
=
+
,
lim
lim
x
u
x
u
x
\to \infty 
\to 
+
=
=
+
. 
\textcolor{lyred}{例}.
(
)(
)
(
)(
)
lim
lim
lim
u
x
x
u
u
u
u
x
x
u
u
u
u
u
u
x
x
=
\to 
\to 
\to 
-
+
-
-
-
-
=
=
=
+
-
-
+
+
-
. 
\textcolor{lyred}{定理}.设复合函数
()
f
g x



在
(
)
U x
\circ 
内有定义,
()
lim
x
x g x
u
\to 
=
,并且,在
(
)
U x
\circ 
内 
()
g x
u
\neq 
,
()
lim
u
u f u
A
\to 
=
,则
()
()
()
lim
lim
u g x
x
x
u
u
f
g x
f u
A
=
\to 
\to 
=
=




. 
\vspace{4pt}
{\bfseries \textcolor{lyred}{三.幂指函数的极限运算法则}}

\textcolor{lyred}{定理}.设
()
lim
x
f x
A
\to 
=
>
\Omega 
,
()
lim
x
g x
B
\to 
=
\Omega 
,则
()
()
lim
g x
B
x
f x
A
\to 
=




\Omega 
. 
\vspace{4pt}
{\bfseries \textcolor{lyred}{四.斜渐近线}}

\textcolor{lyred}{定义}.若
()(
)
lim
x
f x
ax
b
\to +\infty 
-
+
=




,或
()(
)
lim
x
f x
ax
b
\to -\infty 
-
+
=




,则称y
ax
b
=
+
为 
曲线
()
y
f x
=
的\textcolor{lyred}{斜渐近线}. 
\textcolor{lyred}{命题}.
()(
)
()
lim
lim
x
x
f x
f x
ax
b
a
x
\to \infty 
\to \infty 
-
+
=
\Leftrightarrow 
=




,
()
lim
x
f x
ax
b
\to \infty 
-
=




. 
\textcolor{lyred}{例}.求曲线
y
x
x
x
=
-
-
的斜渐近线. 
解.(1)
lim
lim
lim
x
x
x
x
x
x
x
x
a
x
x
x
\to +\infty 
\to +\infty 
\to +\infty 
-
-
-
=
=
-
=
-
-
=, 
(
)
(
)
lim 2
lim
lim
x
x
x
x
b
x
x
x
x
x
x
x
x
x
x
\to +\infty 
\to +\infty 
\to +\infty 
=
-
-
-
=
-
-
=
=
+
-
,故 
y
x
=
+
为斜渐近线; 
(2)
lim
lim
lim
x
x
x
x
x
x
x
x
a
x
x
x
\to -\infty 
\to -\infty 
\to -\infty 
-
-
-
=
=
-
=
+
-
=, 
(
)
(
)
lim 2
lim
lim
x
x
x
x
b
x
x
x
x
x
x
x
x
x
x
\to -\infty 
\to -\infty 
\to -\infty 
=
-
-
-
=-
-
+
=
=-
-
-
,故 
y
x
=
-
也为斜渐近线. 
\textcolor{lyred}{补充练习 }
1.
(
)(
)
(
)(
)
8cos
2cos
lim
lim
lim
2cos
cos
x
u
u
u
u
x
x
u
u
x
x
u
u
u
u
\pi 
\to 
\to 
\to 
-
+
-
-
-
-
=
=
=
+
-
+
-
-
+
. 
2.
lim
lim
x
x
x
x
x
x
x
x
x
x
x
x
\to +\infty 
\to +\infty 
+


+
+
-
=
=




+
+
+
. 
3.
(
)
(
)
(
)
(
)(
)
lim
lim
x
x
x
x
x
x
x
x
x
x
x
x
\to \infty 
\to \infty 
+-
-
+
-
+
=
=
+
+
+
\cdots 
+
+
+
. 
4.
+
+
+
lim
lim
lim
2 2
x
x
x
x
x
x
x
x
x
x
x
x
x
x
x
\to 
\to 
\to 
-
-
-
=
\cdots 
\cdots 
=
=
-
+
+
-
-
(
)(
)
+
+
+
lim
lim
lim
2 2
2 2
2 2
x
x
x
x x
x
x
x
x
x
x
\to 
\to 
\to 
-
+
+
-
=
=
-=
-
-
. 
5.设lim
n
n
x
\to \infty 
存在,且
2lim
n
n
n
x
n
n
n
n
x
\to \infty 
=
-
-
+
+
,求lim
n
n
x
\to \infty 
. 
解.设lim
n
n
x
a
\to \infty 
=
,则
(
)
lim
n
a
n
n
n
n
a
\to \infty 
=
-
-
+
+
,故 
(
)
lim
lim
n
n
n
a
n
n
n
n
n
n
n
n
\to \infty 
\to \infty 
=
+
-
-
=
=
+
+
-
. 
6.设
()
f x 为多项式,且
()
lim
x
f x
x
x
\to \infty 
-
=-,
()
lim
x
f x
x
\to 
=
,求
()
f x . 
解.设
()
f x
x
x
ax
b
=
-
+
+
,
()
lim
lim
x
x
f x
b
a
b
x
x
\to 
\to 


=
+
=
\Rightarrow 
=




,
a =
. 
7.设
()
f x 为三次多项式,且
()
()
lim
lim
x
x
f x
f x
x
x
\to 
\to 
=
=
-
-
,求
()
lim
x
f x
x
\to 
-
. 
解.设
()
(
)(
)(
)
f x
a x
x
x
b
=
-
-
-
,则
()
(
)
lim
x
f x
a
b
x
\to 
=-
-
=
-
, 
()
(
)
lim
x
f x
a
b
x
\to 
=
-
=
-
,解得
a =
,
b =
,故
()
lim
x
f x
x
\to 
=-
-
. 
\vspace{6pt}
{\large \bfseries \textcolor{lyred}{第1.6 节 极限存在准则 两个重要极限}}

\vspace{4pt}
{\bfseries \textcolor{lyred}{一.第一个重要极限}}

\textcolor{lyred}{定理}.(1)设
n
n
n
y
x
z
\le 
\le 
,若lim
lim
n
n
n
n
y
z
a
\to \infty 
\to \infty 
=
=
,则lim
n
n
x
a
\to \infty 
=
; 
(2)设()
()
()
g x
f x
h x
\le 
\le 
,若
()
()
lim
lim
x
x
g x
h x
A
\to 
\to 
=
=
\Omega 
\Omega 
,则
()
lim
x
f x
A
\to 
=
\Omega 
. 
\textcolor{lyred}{例}.设
nx
n
n
n
n
=
+
+
+
+
+
+
\Lambda 
,求lim
n
n
x
\to \infty 
. 
解.
n
n
x
n
n
n
n
\cdots 
\le 
\le 
\cdots 
+
+
,故lim
n
n
x
\to \infty 
=. 
\textcolor{lyred}{例}.设
n
n
n
n
n
m
x
a
a
a
=
+
+
+
\Lambda 
,其中
ia \ge 
,求lim
n
n
x
\to \infty 
. 
解.记
m
a
a
\le 
\le 
\Lambda 
,则
n
n
n
n
m
n
m
a
x
m a
\le 
\le 
\cdots 
,故lim
n
m
n
x
a
\to \infty 
=
. 
\textcolor{lyred}{重要极限1}.
sin
lim
x
x
x
\to 
=. 
\textcolor{lyred}{例}.(1)
tan
sin
sin
lim
lim
lim
lim
cos
cos
x
x
x
x
x
x
x
x
x
x
x
x
\to 
\to 
\to 
\to 


=
=
\cdots 
=




. 
(2)
arcsin
arcsin
lim
lim
sin
t
x
x
t
x
t
x
t
=
\to 
\to 
=
=,
arctan
arctan
lim
lim
tan
t
x
x
t
x
t
x
t
=
\to 
\to 
=
=. 
(3)
tan3
tan3
lim
lim
sin 2
sin 2
x
x
x
x
x
x
x
x
x
x
\to 
\to 


=
\cdots 
\cdots 
=




. 
(4)
(
)
1 cos
2sin
2sin
sin
lim
lim
lim
lim
x
x
x
x
x
x
x
x
x
x
x
x
\to 
\to 
\to 
\to 


-
=
=
=
=




. 
\textcolor{lyred}{例}.(1)
(
)
tan
1 cos
tan
sin
lim
tan
sin
lim
lim
t
x
x
t
t
t
t
t
t
x
x
x
t
t
=
\to \infty 
\to 
\to 
\cdots 
-
-


-
=
=
=




. 
(2)
(
)
(
)
lim
lim
lim
sin
sin
sin
t x
x
t
t
x
t
t
x
t
t
\pi 
\pi 
\pi 
\pi 
\pi 
\pi 
\pi 
=-
\to 
\to 
\to 
-
=
=-
=-
+
. 
(3)
(
)
(
)
(
)
(
)
(
)
sin
sin
sin
lim
lim
lim
sin
sin
sin
m
t x
m n
n
x
t
t
mt
m
mt
mx
m
nx
nt
n
n
nt
\pi 
\pi 
\pi 
\pi 
=-
-
\to 
\to 
\to 
+
-
=
=
=-
+
-
. 
\vspace{4pt}
{\bfseries \textcolor{lyred}{二.第二个重要极限}}

\textcolor{lyred}{定理}.单调有界数列必收敛. 
\textcolor{lyred}{例}.求
(
)
lim
!
n
n
a
a
n
\to \infty 
>
. 
解.
!
n
n
a
x
n
=
,则
n
n
a
x
x
n
+=
\cdots 
+,当
n
a
>
-时,
n
n
x
x
+<
,单减,又
nx >
,故有界, 
因此收敛,设lim
n
n
x
x
\to \infty 
=
,则
n
n
a
x
x
x
x
n
+=
\cdots 
\Rightarrow 
=
\cdots 
=
+
. 
\textcolor{lyred}{例}.设
x =
,
x =
+
,\Lambda ,
n
n
x
x -
=
+
,求lim
n
n
x
\to \infty 
. 
解.显然
x
x
<
,而
n
n
n
n
n
n
x
x
x
x
x
x
-
-
+
<
\Rightarrow 
=
+
<
=
+
,由归纳法,\{\}
nx
单增; 
显然
x <
,而
n
n
n
x
x
x
+
<
\Rightarrow 
=
+
<
,由归纳法,
nx <
; 
设lim
n
n
x
x
\to \infty 
=
,则
x
x
=
+
,故
x =
. 
\textcolor{lyred}{定理}.设
n
nx
n


=
+




,则\{\}
nx
单调有界,因此收敛.记lim
2.71828
n
n
x
e
\to \infty 
=
=
\Lambda . 
\textcolor{lyred}{重要极限2}.
lim 1
x
x
e
x
\to \infty 


+
=




;
(
)
lim 1
x
x
x
e
\to 
+
=. 
\textcolor{lyred}{例}.
(
)
sin
lim 1 sin
x
x
x
e
\to 
+
=
,
(
)
lim 1
x
x
x
e
+
-
\to 
+
-
=
. 
\textcolor{lyred}{例}.(1)
(
)
(
)
1 2
lim 1 2
lim 1 2
x
x
x
x
x
x
e
\cdots 
\to 
\to 
+
=
+
=
. 
(2)
(
)
lim 1
lim 1
x
x
x
x
e
x
x
\cdots -
-
-
\to \infty 
\to \infty 
-




-
=
+
=








. 
\textcolor{lyred}{例}.(1)
(
)
1 1
lim
lim 1
x
x
x
x
x
x
x
x
e
-
\cdots 
-
-
-
-
\to 
\to 
=
+
-
=
. 
(2)
1 2
1 2
lim
lim 1
lim 1
1 2
1 2
1 2
x
x
x
x
x
x x
x
x
x
x
x
x
e
x
x
x
-
\cdots 
\cdots 
-
\to 
\to 
\to 
+
+






=
+
-
=
+
=






-
-
-






. 
(3)
1 3
1 3
1 2
1 3
lim
lim 1
n
n
n
n
n
n
n
n
n
n
n
n
e
n
n
n
n
-
-
\cdots 
\cdots 
-
-
\to \infty 
\to \infty 


-
-


=
+
=




-
-




. 
(4)
(
)
(
)
cos
cos
lim cos
lim 1
cos
x
x
x
x
x
x
x
x
e
-
-
\cdots 
-
\to 
\to 
=
+
-
=




. 
(5)
sin
cos
sin
cos
lim sin
cos
lim 1 sin
cos
x
x
x
x
x
x
x
x
e
x
x
x
x
+
-
\cdots 
+
-
\to \infty 
\to \infty 




+
=
+
+
-
=








. 
\textcolor{lyred}{例}.(1)
(
)
(
)
ln 1
lim
limln 1
limln
ln
x
x
x
u
e
x
x
u
e
x
\to 
\to 
\to 
+
=
+
=
=
=. 
(2)
(
)(
)
(
)
ln 1
ln
lim
lim
x
x
x
x
x
x
x
x
\to 
\to 
+
-
=
\cdots 
=
+
-
+
-
. 
\textcolor{lyred}{例}.(1)
(
)
lim
lim
ln 1
x
x
t e
x
t
e
t
x
t
=
-
\to 
\to 
-
=
=
+
. 
(2)
cos
cos
cos
1 cos
lim
lim
lim cos
x
x
x
x
x
x
e
e
e
e
x
e
e
e
x
x
x
x
-
-
\to 
\to 
\to 
-
-
-
-
=\cdots 
=\cdots 
\cdots 
=-
-
. 
(3)
cos
1 cos
lim
lim
lim
x
x
x
x
x
e
x
e
x
x
x
x
\to 
\to 
\to 
-
-
-
=
+
=+
=
. 
\textcolor{lyred}{补充练习 }
1.证明:(1)
!
lim
n
n
n
n
\to \infty 
=
;(2)
lim 1
n
n
n
\to \infty 
+
+
+
=
\Lambda 
. 
证.(1)
!
n
n
n
n
<
<
;(2)
n
n n
n
<
+
+
+
<
\Lambda 
,证毕. 
2.
sin 2
tan5
sin 2
tan5
lim
lim
sin3
tan 2
sin3
tan 2
x
x
x
x
x
x
x
x
x
x
x
x
x
x
\to 
\to 
+
+
+
=
=
=
-
-
-
. 
3.
(
)
cos2
cos2
cos2
cos 2
1 cos2
lim
lim
lim
cos2
cos2
x
x
x
x
x
x
x
x
x
x
x
x
x
\to 
\to 
\to 
-
-
-
=
=
=
+
. 
4.设
()
lim
x
f x
\pi 
\to 
存在,且
()
()
sin
2lim
x
x
f x
f x
x
\pi 
\pi 
\to 
=
+
-
,求
()
f x . 
解.设
()
lim
x
f x
A
\pi 
\to 
=
,则
()
sin
x
f
x
A
x
\pi 
=
+
-
,取极限,得
sin
lim
x
x
A
A
x
\pi 
\pi 
\to 
=
+
-
, 
(
)
(
)
sin
sin
sin
lim
lim
lim
x
x
x
x
x
x
A
x
x
x
\pi 
\pi 
\pi 
\pi 
\pi 
\pi 
\pi 
\pi 
\pi 
\to 
\to 
\to 
-
+
-
=-
=-
=
=
-
-
-
,
()
sin
x
f x
x
\pi 
=
+
-
. 
5.讨论
sin
lim
x
x
x
e
x
x
e
\to 


+


+




+


的存在性. 
解.
sin
lim
2 1
x
x
x
e
x
x
e
-
\to 


+


+
=
-=




+


,
sin
lim
0 1
x
x
x
e
x
x
e
+
\to 


+


+
=
+=




+


,故存在. 
6.
(
)
lim
sin
lim 1
sin
n
n
n
n
n
n
n
e
n
n
n
n
+
+
\to \infty 
\to \infty 
+


=
+
\cdots 
=




. 
7.
1 sin
tan
sin
tan
sin
1 sin
tan
tan
sin
lim
lim 1
1 sin
1 sin
x
x
x
x
x
x
x
x
x
x
x
x
x
e
x
x
+
-
\cdots 
\cdots 
-
+
\to 
\to 
+
-




=
+
=




+
+




. 
8.
(
)
(
)
cos
cos
lim cos
lim 1
cos
x
x
x
x
x
x
e
e
x
x
x
x
e
-
-
-
\cdots 
\cdots 
-
-
-
\to 
\to 
=
+
-
=




. 
9.
sin
sin
cos
cos
lim
lim
lim
x
x
x
x
x
x
e
x
e
x
x
x
\to 
\to 
\to 
-
-
-
=
+
=-
-=-
. 
10.
(
)
lim
lim
lim
n n
n
n
n
n
n
n
n
n
n
e
e
e
n
e
n
e
-
+
+
+
+
\to \infty 
\to \infty 
\to \infty 






-
=
\cdots 
-
=
-
=














(
)
(
)
(
)
lim
n n
n
e
n
n n
n n
+
\to \infty 
-
\cdots 
\cdots 
=
+
+
. 
\vspace{6pt}
{\large \bfseries \textcolor{lyred}{第1.7 节 无穷小的比较}}

\vspace{4pt}
{\bfseries \textcolor{lyred}{一.无穷小的阶}}

\textcolor{lyred}{定义}.设\alpha , \beta 为x \to \Omega 过程中的无穷小量,则 
(1)当lim
x
\beta 
\alpha 
\to 
=
\Omega 
时,称\beta 是比\alpha \textcolor{lyred}{高阶的无穷小},记为
()
o
\beta 
\alpha 
=
; 
(2)当lim
x
\beta 
\alpha 
\to 
=\infty 
\Omega 
时,称\beta 是比\alpha \textcolor{lyred}{低阶的无穷小},记为
()
o
\alpha 
\beta 
=
; 
(3)当lim
x
C
\beta 
\alpha 
\to 
=
\neq 
\Omega 
时,称为\textcolor{lyred}{同阶无穷小};若lim
x
\beta 
\alpha 
\to 
=
\Omega 
,则称为\textcolor{lyred}{等价},记
~
\beta 
\alpha ; 
\textcolor{lyred}{注}.等价关系具有\textcolor{lyred}{传递性},即
~
\alpha 
\beta ,
~
~
\beta 
\gamma 
\alpha 
\gamma 
\Rightarrow 
. 
\textcolor{lyred}{例}.当
x \to 
时,sin
~
x
x ,arcsin
~
x
x , tan
~
x
x ,arctan
~
x
x ,
(
)
ln 1
~
x
x
+
,\textcolor{lyred}{ }
1 ~
xe
x
-
,
ln
1~
ln
x
x
a
a
e
x
a
-=
-
,(
)
(
)
(
)
ln 1
1~
ln 1
~
x
x
e
x
x
\alpha 
\alpha 
\alpha 
\alpha 
+
+
-=
-
+
, 
1~
n
x
x
n
+
-
,
cos
~ 2
x
x
-
. 
\vspace{4pt}
{\bfseries \textcolor{lyred}{二.近似公式}}

\textcolor{lyred}{定理}.当x \to \Omega 时,
()
~
o
\beta 
\alpha 
\beta 
\alpha 
\alpha 
\Leftrightarrow 
=
+
,此时称\alpha 为\beta 的\textcolor{lyred}{主部}. 
\textcolor{lyred}{注}.
()
sin x
x
o x
=
+
,
()
tan x
x
o x
=
+
,
(
)
()
ln 1
x
x
o x
+
=
+
,
()
xe
x
o x
=+
+
, 
(
)
()
x
x
o x
\alpha 
\alpha 
+
=+
+
,
(
)
cos
x
x
o x
=-
+
. 
\textcolor{lyred}{例}.
(
)
()
()
()
()
()
lim
lim
lim
ln 1
x
x
x
x
x
o x
x
o x
x
o x
x
e
x
x
o x
x
o x
\to 
\to 
\to 




+
+
-
+
+
+




+
-




=
=
=
+
+
+
. 
\textcolor{lyred}{例}.
()
()
()
()
lim
lim
x
x
x
o x
x
o x
x
x
x
x
x
o x
x
o x
\to 
\to 




+
+
-
-
+




+
-
-




=
=




+
-
-
+
+
-
-
+








. 
\vspace{4pt}
{\bfseries \textcolor{lyred}{三.等价无穷小的替换法则}}

\textcolor{lyred}{定理}.设当x \to \Omega 时,
~
\alpha 
\alpha ',
~
\beta 
\beta ',则lim
lim
x
x
\beta 
\beta 
\alpha 
\alpha 
\to 
\to 
'
=
'
\Omega 
\Omega 
. 
\textcolor{lyred}{注(因式替代原则)}.若当x \to \Omega 时,
~
\alpha 
\alpha ',则
()
()
lim
lim
x
x
f x
f x
\alpha 
\alpha 
\to 
\to 
'
=
\Omega 
\Omega 
. 
\textcolor{lyred}{注}.经典的错误:
tan
sin
lim
lim
sin
x
x
x
x
x
x
x
x
\to 
\to 
-
-
=
=
; 
tan
sin
tan
sin
lim
lim
lim
lim
lim
lim
sin
sin
sin
x
x
x
x
x
x
x
x
x
x
x
x
x
x
x
x
x
x
x
x
\to 
\to 
\to 
\to 
\to 
\to 
-
-
=
-
=
-
=
=
, 
实际上:
(
)
tan
1 cos
tan
sin
lim
lim
lim
sin
sin
x
x
x
x
x
x
x
x
x
x
x
x
\to 
\to 
\to 
-
\cdots 
-
=
=
=
. 
\textcolor{lyred}{例}.(1)
sin5
lim
lim
x
x
x
x
x
x
x
\to 
\to 
=
=
+
, lim 2 sin
lim 2
n
n
n
n
n
n
x
x
x
\to \infty 
\to \infty 
=
\cdots 
=
. 
(2)
lim sin ln 1
lim ln 1
lim
x
x
x
x
x
x
x
x
x
\to \infty 
\to \infty 
\to \infty 




+
=
+
=
\cdots 
=








. 
(3)
(
)
(
)
(
)
(
)
1 cos
lim
lim
lim
ln 1 3
x
x
x
x
x
e
e
x
x
x
x
\to 
\to 
\to 
-
-
-
=
=
=-
-
-
-
. 
(4)
(
)
cos
cos
1 1
cos
lim
lim
lim
ln 1
tan
tan
x
x
x
x
x
x
x
x
x
x
x
x
x
x
\to 
\to 
\to 
+
-
+
+
--
+
-
=
=
=
-
-
-
. 
\textcolor{lyred}{例}.(1)
tan
lim tan3 tan
lim tan3
tan
lim
tan3
t
x
t
t
x
t
x
x
t
t
t
\pi 
\pi 
\pi 
\pi 
=
-
\to 
\to 
\to 




-
=
-
=
=








. 
(2)
(
)
1 8
lim
lim
lim
t
x
x
t
t
t
t
x
x
x
t
t
t
t
=
\to \infty 
\to 
\to 


+
-
+
+
-
+
=
+
-
+
=
=






1 8
lim
lim
t
t
t
t
t
t
\to 
\to 
+
-
+-
-
=
-
=
. 
\textcolor{lyred}{例}.(1)
tan
sin
tan
sin
sin
tan
sin
lim
lim
lim
sin
sin
x
x
x
x
x
x
x
x
e
e
e
x
x
e
x
x
x
-
\to 
\to 
\to 
-
-
-
=
\cdots 
=
=
. 
(2)
(
)
ln 1 3
ln 1
lim
lim3
lim
lim
x
x
x
x
x
x
x
x
x
x
x
x
x
e
x
x
x
x


+




\to 
\to 
\to 
\to 




+
-
+




+
-
-




=
=
=
=
. 
(3)
ln
ln
ln
lim
lim
lim
lim
x
x
x
x
x
x
x
x
a
b
a
b
a
b
a
b
x
x
x
x
x
x
x
x
a
b
e
e
e
e
ab
\to 
\to 




+
+
+
-
+
-












\to 
\to 


+
=
=
=
=
=




. 
\textcolor{lyred}{注}.
lim
x
x
x
k
k
k
x
a
a
a
a
k
\to 


+
+
=




\Lambda 
\Lambda 
,
lim
n
n
n
k
k
k
n
a
a
a
a
k
\to \infty 


+
+
=






\Lambda 
\Lambda 
. 
(4)
lim
ln sin
lim
sin
sin
lim
1 ln2
lim sin
t
x
x
x
x
t
x
x
x
t
x
x
x
t
x
e
e
e
e
e
x
\to \infty 
\to \infty 
\to 








+
+
-
+
-




+




\to \infty 


+
=
=
=
=
=




. 
\textcolor{lyred}{例}.已知
()
lim
x
f x
x
\to 
=
,求
()
sin
lim 1
x
x
f x
x
\to 


+




. 
解.
()
lim
x
f x
x
\to 
=
,故
()
()
()
ln 1
sin
lim
sin
lim 1
lim
x
f x
f x
x
x
x
x
x
x
f x
e
e
e
x
\to 


+




\to 
\to 


+
=
=
=




; 
或者,
()
()
()
()
()
sin
sin
lim
lim 1
lim 1
x
f x
x
f x
x
f x
x
x
x
x
x
f x
f x
e
e
x
x
\to 
\cdots 
\cdots 
\to 
\to 




+
=
+
=
=








. 
\textcolor{lyred}{补充练习 }
1.
lim
sin
lim
lim
x
x
x
x
x
x
x
x
x
x
x
\to -\infty 
\to -\infty 
\to -\infty 
=
\cdots 
=
=
=-
-
-
-
-
. 
2.
lim
lim
lim
cos
cos
cos
x
x
x
x
x
x
x
x
x
x
\to 
\to 
\to 
+
-
+
+
-
+
-
=
-
=
-
=-
-
-
-
-
-
. 
3.
(
)
(
)(
)
(
)
2 1
lim
lim
lim
ln
ln 1
x
x
x
x
x
e
e
e
x
e
e
e
x x
x
x
x
x
x
x
-
\to 
\to 
\to 
-
-
-
=
=
=
-
+
-
+
-
+
-
. 
4.
1 sin
1 sin
1 sin
1 sin
lim
lim
lim
lim
t
t
t
x
x
x
t
t
t
e
e
e
t
x e
e
e
e
t
t
t
=
-
-
-
-
\to \infty 
\to 
\to 
\to 


-
-
-
-
-
=
=
=
=






sin
lim
t
t
e
e
t
\to 
-
=-
. 
5.
(
)
ln 2
lim
lim
lim
ln 2
x x
x
x
x
x
x
x
x
x
x
x
x
e
-
+
+
+
+
\to \infty 
\to \infty 
\to \infty 






-
=
-
=
-
=














. 
6.
(
)
2 cos
ln
cos
lim
lim
ln cos2
cos2
x
x
x
x
x
x
e
x
x
x
x
+






\to 
\to 




+

-
=
-
=






-










cos
cos
1 cos
lim
ln
lim
ln 1
lim
x
x
x
x
x
x
x
x
x
\to 
\to 
\to 
-
+
-
-
-
-




=
+
=
\cdots 
=








. 
\vspace{6pt}
{\large \bfseries \textcolor{lyred}{第1.8 节 函数的连续性与间断点}}

\vspace{4pt}
{\bfseries \textcolor{lyred}{一.函数的连续性}}

设
()
y
f x
=
在
(
)
U x
内有定义,记
x
x
x
\Delta =
-
,称为x 的\textcolor{lyred}{增量},相应地,有函数值的 
\textcolor{lyred}{增量}
()
(
)
(
)
(
)
y
f x
f x
f x
x
f x
\Delta 
\Delta 
=
-
=
+
-
. 
\textcolor{lyred}{定义}.若
lim
x
y
\Delta 
\Delta 
\to 
=
,或
()
(
)
lim
x
x f x
f x
\to 
=
,则称
()
y
f x
=
在
0x 处\textcolor{lyred}{连续}. 
\textcolor{lyred}{定理}.
()
f x 在
0x 处连续
(
)
(
)
(
)
f x
f x
f x
-
+
\Leftrightarrow 
=
=
. 
\textcolor{lyred}{定义}.设
()
f x 在
0x 的左侧邻域(
)
,
x
x
\delta 
-
内有定义,若(
)
(
)
f x
f x
-=
,则称
()
f x 
在
0x 处\textcolor{lyred}{左连续};类似地,可以定义\textcolor{lyred}{右连续}. 
\textcolor{lyred}{定义}.若
()
f x 在(
)
,a b 内处处连续,则称
()
f x 为(
)
,a b 上的\textcolor{lyred}{连续函数}; 
若
()
f x 在(
)
,a b 内处处连续,且在x
a
=
处右连续,在x
b
=
处左连续,则称
()
f x 为 
闭区间[
]
,a b 上的\textcolor{lyred}{连续函数}. 
\textcolor{lyred}{例}.设
()
(
)
cos2 ,
ln 1
,
x
x
f x
ax
x
x
\le 

=
+

>

,若
()
f x 在
x =
处连续,求a . 
解. (
)
lim cos2
x
f
x
-
-
\to 
=
=, (
)
(
)
ln 1
lim
x
ax
f
a
x
+
+
\to 
+
=
=
,
()
f
=,故
a =. 
\vspace{4pt}
{\bfseries \textcolor{lyred}{二.函数的间断点}}

\textcolor{lyred}{定义}.设
()
f x 在
(
)
U x
\circ 
内有定义,且在
0x 处不连续,即有以下三种情况之一发生: 
(1)
(
)
f x
没有定义;(2)极限
()
lim
x
x f x
\to 
不存在;(3)
()
(
)
lim
x
x f x
f x
\to 
\neq 
,则称
0x 为 
()
f x 的\textcolor{lyred}{间断点}. 
\textcolor{lyred}{第一类间断点}: (
)
f x-与(
)
f x+均存在. 
(1)若(
)
(
)
f x
f x
-
+
\neq 
,则称
0x 为\textcolor{lyred}{跳跃间断点},例如sgn x 在
x =
处. 
(2)若(
)
(
)
f x
f x
-
+
=
,即
()
lim
x
x f x
\to 
存在,则称
0x 为\textcolor{lyred}{可去间断点}. 
\textcolor{lyred}{第二类间断点}: (
)
f x-与(
)
f x+中至少有一个不存在. 
特别地,若(
)
f x-或(
)
f x+=\infty ,则称
0x 为\textcolor{lyred}{无穷间断点},例如tan x 在
x
\pi 
=\pm 
处. 
\textcolor{lyred}{例}.讨论
()
sin x
f x
x
=
在
x =
处的间断点类型. 
解.
()
f x 在
x =
处没有定义,而
()
lim
x
f x
\to 
=存在,故为可去间断点; 
若补充定义
()
f
=,则
x =
成为连续点. 
\textcolor{lyred}{例}.设
()
1,
1,
x
x
f x
x
x
-
\le 

=+
>

,讨论
()
f x 在
x =
处的间断点类型. 
解. (
)
(
)
lim
x
f
x
-
-
\to 
=
-
=-, (
)
(
)
lim
x
f
x
+
+
\to 
=
+
=,故为跳跃间断点. 
\textcolor{lyred}{例}.讨论
()
arctan
f x
x
=
在
x =
处的间断点类型. 
解. (
)
f
\pi 
-=-
, (
)
f
\pi 
+=
,故
x =
为跳跃间断点. 
\textcolor{lyred}{例}.讨论
()
sin
f x
x
=
在
x =
处的间断点类型. 
解. (
)
f
-与(
)
f
+均不存在,故
x =
为第二类间断点. 
\textcolor{lyred}{例}.确定
()
x
f x
x
x
-
=
-
+
的间断点及其类型.\textcolor{lyred}{ }
解.
()
(
)(
)
(
)(
)
x
x
f x
x
x
-
+
=
-
-
,故
x =为可去间断点,
x =
为无穷间断点. 
\textcolor{lyred}{例}.确定
()
tan
f x
x
=+
的间断点及其类型. 
解.
x
k
\pi 
\pi 
=
+
为可去间断点,
x
k
\pi 
\pi 
=
-
为无穷间断点. 
\textcolor{lyred}{例}.确定
()
x
x
f x
-
=
+
的间断点及其类型. 
解.
lim 2
x
x
-
\to 
=
,
()
lim 2
lim
x
x
x
f
x
+
-
\to 
\to 
=+\infty \Rightarrow 
=,
()
lim
x
f
x
+
\to 
=-,故
x =
为 
跳跃间断点. 
\textcolor{lyred}{例}.确定
()
lim1
n
n
x
f x
x
\to \infty 
+
=
+
的间断点及其类型.\textcolor{lyred}{ }
解.
()
,
0,
1,
0,
x
x
x
f x
x
x
+
<

>

=
=


=-

,故
x =为跳跃间断点,
x =-为连续点. 
\textcolor{lyred}{补充练习 }
1.确定
()
x
x
f x
e -
=
-
的间断点及其类型.\textcolor{lyred}{ }
解.
()
lim
lim
x
x
x
x
e
e
f x
-
\to 
\to 
=
=\Rightarrow 
=\infty ,故
x =
为无穷间断点; 
()
lim
lim
x
x
x
x
e
f
x
-
-
-
\to 
\to 
=
\Rightarrow 
=,
()
lim
lim
x
x
x
x
e
f
x
+
+
-
\to 
\to 
=+\infty \Rightarrow 
=
,故
x =为跳跃间断点. 
2.确定
()
ln
x x
f x
x
+
=
的间断点及其类型. 
解.
x =为无穷间断点,
x =
为可去间断点; 
(
)
(
)
(
)
(
)
lim
lim
lim
lim
ln
ln
ln
ln 1
t x
x
x
x
t
x x
x x
x
t
x
x
x
t
-
-
-
-
=+
\to -
\to -
\to -
\to 
+
-
+
+
=
=
=
=-
-
-
-
, 
(
)
(
)
lim
lim
ln
ln
x
x
x x
x x
x
x
+
+
\to -
\to -
+
+
=
=
-
,故
x =-为跳跃间断点. 
\vspace{6pt}
{\large \bfseries \textcolor{lyred}{第1.9 节 连续函数的运算与初等函数的连续性}}

\vspace{4pt}
{\bfseries \textcolor{lyred}{一.连续函数的四则运算}}

\textcolor{lyred}{定理}.设
()
f x 和()
g x 均在
0x 处连续,则它们的和差积商(分母不为零)也在
0x 处 
连续. 
\textcolor{lyred}{注}.设在
0x 处
()
f x 连续,而()
g x 间断,则
()
()
f x
g x
\pm 
必在
0x 处间断. 
\textcolor{lyred}{推论}.连续函数的和差积商在定义区间内均是连续函数. 
\vspace{4pt}
{\bfseries \textcolor{lyred}{二.连续函数的反函数与复合函数}}

\textcolor{lyred}{定理}.设
()
y
f x
=
在区间I 上单调且连续,则其反函数
()
x
f
y
-
=
在对应的区间 
f
J
R
=
上单调且连续. 
\textcolor{lyred}{定理}.设
()
lim
x
g x
u
\to 
=
\Omega 
,
()
f u 在
0u 处连续,则
()
(
)
()
lim
lim
x
x
f
g x
f u
f
g x
\to 
\to 


=
=






\Omega 
\Omega 
. 
\textcolor{lyred}{例}.
lim
lim
x
x
x
x
x
x
\to 
\to 
-
-
=
=
-
-
,
sin
sin
lim 2
lim
x
x
x
x
x
x
e
e
e
\to 


-
-




\to 
=
=
. 
\textcolor{lyred}{例}.
(
)
(
)
(
)
log
lim
lim log
log
lim 1
log
ln
a
x
x
a
a
a
x
x
x
x
x
x
e
x
a
\to 
\to 
\to 
+




=
+
=
+
=
=








. 
\textcolor{lyred}{定理}.设
()
u
g x
=
在
x
x
=
处连续,
()
y
f u
=
在
(
)
u
g x
=
处连续,则
()
y
f
g x
=



 
在
x
x
=
处连续. 
\textcolor{lyred}{推论}.连续函数的复合函数在定义区间内连续. 
\textcolor{lyred}{例}.证明:
()
f x 连续
()
f x
\Rightarrow 
也连续. 
证.
()
y
f x
=
由y
u
=
和
()
u
f x
=
复合而成,两者均连续,即得,证毕. 
\textcolor{lyred}{例}.设
()
f x , ()
g x 连续,证明
()
()
\{
\}
max
,
f x
g x
,
()
()
\{
\}
min
,
f x
g x
连续. 
证.
()
()
\{
\}
()
()
()
()
max
,
f x
g x
f x
g x
f x
g x
+
+
-
=
, 
()
()
\{
\}
()
()
()
()
min
,
f x
g x
f x
g x
f x
g x
+
-
-
=
,即得,证毕. 
\vspace{4pt}
{\bfseries \textcolor{lyred}{三.初等函数的连续性}}

\textcolor{lyred}{定理}.所有初等函数在其定义区间内均为连续函数. 
\textcolor{lyred}{例}.
limln arcsin
ln arcsin
ln
x
x
\pi 
\to 
=
=
. 
\textcolor{lyred}{例}.
lim
lim
lim
1 8
x
x
x
x
x
x
x
x
x
\to 
\to 
\to 
+
-
+
-
+
+
=
=
\cdots 
=
-
-
+
+
+
+
. 
\textcolor{lyred}{例}.
tan
1 sin
tan
sin
1 sin
lim
lim
sin
tan
1 sin
1 sin
x
x
x
x
x
x
x
x
x
x
x
x
x
x
\to 
\to 
+
-
+
-
+
+
=
=
+
+
+
+
-
. 
\textcolor{lyred}{补充练习 }
1.
(
)
(
)
(
)
limsin
lim
sin
lim
sin
n
n
n
n
n
n
n
n
n
n
\pi 
\pi 
\pi 
\to \infty 
\to \infty 
\to \infty 
+
=
-
+-
=
-
=
++
. 
2.设
()
f x 在
x =
处连续,且在(
)
,
-\infty +\infty 上满足
()
(
)
f x
f
x
=
,证明:
()
f x 为 
常数函数. 
证.
x
\forall ,
()
2n
x
x
x
f x
f
f
f






=
=
=
=












\Lambda 
,故
()
()
lim
2n
n
x
f
x
f
f
\to \infty 


=
=




,即得, 
证毕. 
3.设
(
)
()
()
f x
y
f x f
y
+
=
,
(
)
,
,
x y
\forall 
\in -\infty +\infty ,并且
()
f x 在
x =
处连续,证明: 
()
f x 在(
)
,
-\infty +\infty 上连续. 
证.
x
\forall ,
()
()
()
f x
f x f
=
,故(1)
()
()
f
f x
=
\Rightarrow 
\equiv 
;或者 
(2)
()
()
(
)
()
(
)
()
(
)
lim
lim
lim
x
x
x
f
f
f x
x
f x f
x
f x
f
x
\Delta \to 
\Delta \to 
\Delta \to 
\neq 
\Rightarrow 
=\Rightarrow 
+\Delta 
=
\Delta 
=
\Delta 
= 
()
()
()
f x f
f x
=
,即得,证毕. 
\vspace{6pt}
{\large \bfseries \textcolor{lyred}{第1.10 节 闭区间上连续函数的性质}}

\vspace{4pt}
{\bfseries \textcolor{lyred}{一.最大值最小值定理}}

\textcolor{lyred}{定理}.设
()
f x 在[
]
,a b 上连续,则存在
1\psi ,
[
]
,a b
\psi \in 
,使得
[
]
,
x
a b
\forall \in 
,均有 
()
()
(
)
f
f x
f
\psi 
\psi 
\le 
\le 
. 
\textcolor{lyred}{推论}.设
()
f x 在[
]
,a b 上连续,则
()
f x 在[
]
,a b 上有界. 
\textcolor{lyred}{推论}.设
()
f x 在(
)
,a b 上连续,且(
)
f a+与(
)
f b-存在,则
()
f x 在(
)
,a b 上有界. 
\vspace{4pt}
{\bfseries \textcolor{lyred}{二.零点定理与介值定理}}

\textcolor{lyred}{定理}.设
()
f x 在[
]
,a b 上连续,且
()
()
f a
f b
\cdots 
<
,则
(
)
,a b
\psi 
\exists \in 
,使得
()
f \psi =
. 
\textcolor{lyred}{例}.证明:方程
x
x
-
+=
在(
)
0,1 上至少有一个根. 
证.设
()
f x
x
x
=
-
+,则
()
f
=>
,
()1
f
=-<
,即得,证毕. 
\textcolor{lyred}{例}.设
()
[
]
0,1
f x
C
\in 
,且
()
f x
\le 
\le ,证明:
[
]
0,1
\psi 
\exists \in 
,使得
()
f \psi 
\psi 
=
. 
证.令
()
()
F x
f x
x
=
-
,则
()
()
F
f
=
\ge 
,
()
()
F
f
=
-\le 
,即得,证毕. 
\textcolor{lyred}{例}.设
()
[
]
0,1
f x
C
\in 
,且
()
()
f
f
=
,证明:
[
]
0,1
\psi 
\exists \in 
,使得
()
f
f
\psi 
\psi 


+
=




. 
证.令
()
()
F x
f
x
f x


=
+
-




,则
()
F
F 

+
=




,故
()
F
F 

\cdots 
\le 




,证毕. 
\textcolor{lyred}{定理}.设
()
f x 在[
]
,a b 上连续,
()
f a
A
=
,
()
f b
B
=
, A
B
\neq 
,则对介于
,
A B 之间的 
任何数C ,存在
(
)
,a b
\psi \in 
,使得
()
f
C
\psi =
. 
\textcolor{lyred}{推论}.有界闭区间上的连续函数能取到最大最小值之间的任何值. 
\textcolor{lyred}{推论}.设
()
[
]
,
f x
C a b
\in 
,
[
]
,
ix
a b
\in 
,
i\lambda >
,
1,2,
,
i
n
=
\Lambda 
,则
[
]
,a b
\psi 
\exists \in 
,使得 
()
()
(
)
(
)
n
n
n
f x
f x
f x
f
\lambda 
\lambda 
\lambda 
\psi 
\lambda 
\lambda 
\lambda 
+
+
+
=
+
+
+
\Lambda 
\Lambda 
. 
\textcolor{lyred}{补充练习 }
1.设
()
[
]
0,1
f x
C
\in 
,且
()
f
=
,
()1
f
=,证明:
(
)
0,1
\psi 
\exists \in 
,使得
()
1 2
f \psi 
\psi 
=-
. 
证.令
()
()1
F x
f x
x
=
-+
,则
()
F
=-,
()
F
=
,即得,证毕. 
2.设
()
(
)
,
f
x
C
\in 
-\infty +\infty ,且
()
lim
x
f x
x
\to \infty 
=
,证明:方程
()
f x
x
+
=
存在实根. 
证.令
()
()
F x
f x
x
=
+
,则
()
lim
x
F x
x
\to \infty 
=>
,由保号性知,当x
X
>
时,
()
F x
x
>
, 
即当x
X
>
时
()
F x >
,当x
X
<-
时
()
F x <
,即得,证毕. 